\section{Publishing and Querying}
\label{sec:publishing}

\term{python -m vidbpedia serve} starts one FastAPI application that hosts both
the Linked Data routes and a Gradio~\cite{abid2019gradio} user interface. At
start-up it loads the full N-Triples dump into an in-memory rdflib graph,
loads \term{inferred.nt} into a set so that every triple can be labelled as
asserted or inferred, builds a search index over labels and alternative names
that ignores Vietnamese diacritics, and precomputes the class tree. The Linked
Data routes are registered before Gradio is mounted at
\term{/}, because the catch-all route of Gradio would otherwise shadow them.

\subsection{Dereferenceable IRIs}
\label{sec:deref}

The server follows the pattern of DBpedia and of the W3C Cool URIs
note~\cite{sauermann2008cooluris}. Looking up a resource IRI never returns
content directly. It answers \emph{303 See Other} and redirects either to an
HTML page or to an RDF document, depending on the \term{Accept} header
(Table~\ref{tab:routes}). Because the interface is a single-page application,
the HTML page is the address \term{/?resource=<name>}. When this address is
loaded, the interface opens the resource tab on that entity, so every page has
its own URL that can be shared.

\begin{table}[t]
\centering
\caption{Linked Data routes (\term{web/linked\_data.py}). An unknown name returns 404.}
\label{tab:routes}
\small
\begin{tabularx}{\textwidth}{@{}l>{\raggedright\arraybackslash}X@{}}
\toprule
Route & Response \\
\midrule
\term{GET /resource/\{name\}} & 303 to \term{/?resource=\{name\}} for browsers, or to
  \term{/data/\{name\}.ttl}, \term{.nt}, \term{.jsonld} or \term{.rdf} when \term{Accept} asks for
  Turtle, N-Triples, JSON-LD or RDF/XML; sends \term{Vary: Accept} \\
\term{GET /page/\{name\}} & 303 to the HTML page, like \term{dbpedia.org/page/…} \\
\term{GET /data/\{name\}.\{ext\}} & every triple with the resource as subject, and up to 2{,}000 with it
  as object, in the requested format; sends \term{Access-Control-Allow-Origin: *} \\
\term{GET /ontology/\{term\}} & the definition of a \term{vio:} class or property, in Turtle \\
\bottomrule
\end{tabularx}
\end{table}

Listing~\ref{lst:curl} shows an exchange recorded on the running server. As
browsers do, clients must percent-encode the non-ASCII characters of the path.

\begin{listing}[t]
\begin{Verbatim}
$ curl -si -H "Accept: text/turtle" \
       http://127.0.0.1:7860/resource/Nguy%E1%BB%85n_C%C3%B4ng_Ph%C6%B0%E1%BB%A3ng
HTTP/1.1 303 See Other
vary: Accept
location: /data/Nguy%E1%BB%85n_C%C3%B4ng_Ph%C6%B0%E1%BB%A3ng.ttl

$ curl -sL -H "Accept: text/turtle" http://127.0.0.1:7860/resource/Nguy%E1%BB%85n_…
vres:Mito_HollyHock vio:hasPlayer vres:Nguyễn_Công_Phượng .
…
vres:Nguyễn_Công_Phượng a dbo:Animal, dbo:Athlete, dbo:Person, dbo:SoccerPlayer, … ;
    rdfs:label "Nguyễn Công Phượng"@en, "Nguyễn Công Phượng"@vi ;
    dbo:birthDate "1995-01-21"^^xsd:date ;
    …
    vio:birthProvince vres:Nghệ_An ;
    vio:careerStation vres:Nguyễn_Công_Phượng__1, … ;
    …
    owl:sameAs dbr:Nguyễn_Công_Phượng, wd:Q18045362 ;
\end{Verbatim}
\caption{Dereferencing a resource IRI with \term{curl}. Response headers are
abridged and the Turtle is an excerpt.}
\label{lst:curl}
\end{listing}

In the response of Listing~\ref{lst:curl}, the types \term{dbo:*} and the edge \term{vio:hasPlayer} are all inferred;
\term{dbo:Animal} appears because the DBpedia Ontology places \term{dbo:Person}
under \term{dbo:Animal}, and our alignment module (\term{050-dbo-alignment.ttl})
copies the DBpedia class chain.

\subsection{Web Interface}
\label{sec:ui}

The interface has four tabs, shown in Figure~\ref{fig:ui}.
\begin{itemize}
  \item \textbf{Resource tree} (the default tab) lists all 8{,}624 resources
        under the \term{vio:} class tree. Each resource appears under its most
        specific class only, like a file under its folder. Categories,
        redirects and resources that have only a label are listed in separate
        groups, and a filter box matches names with or without diacritics.
  \item \textbf{Resource} renders a page modelled on
        \term{dbpedia.org/page/…}. It shows the abstract, the image, and
        links to Wikipedia, DBpedia, Wikidata and OpenStreetMap. Below come
        the class tree with an asserted or inferred badge for each class, the
        Linked Data links with RDF downloads in four formats, an SVG
        neighbourhood graph with inferred edges drawn dashed, and a
        collapsible relation tree up to three hops deep (player
        $\rightarrow$ club $\rightarrow$ stadium $\rightarrow$ province). The
        page ends with the full property table and the incoming references.
  \item \textbf{Question answering} (Section~\ref{sec:qa}).
  \item \textbf{SPARQL} is an editor where the 19 project prefixes are
        predeclared, with nine example queries. It shows the result as a
        table, or as SPARQL~1.1 JSON or CSV results.
\end{itemize}

\begin{figure}[t]
\centering
\begin{subfigure}[t]{8.55cm}
  \centering
  \includegraphics[height=7cm]{ui_resource.png}
  \caption{Resource page (abstract and image hidden): class tree with
  asserted (\emph{khai báo}) and inferred (\emph{suy luận}) classes, and the
  Linked Data card.}
\end{subfigure}\hfill
\begin{subfigure}[t]{5.95cm}
  \centering
  \includegraphics[height=7cm]{ui_sparql.png}
  \caption{SPARQL tab: counting the instances of DBpedia classes, all of
  them inferred.}
\end{subfigure}
\caption{Web interface. Its labels are in Vietnamese.}
\label{fig:ui}
\end{figure}

\subsection{Natural-Language Question Answering}
\label{sec:qa}

The question-answering tab lets users ask questions in Vietnamese. A language
model writes a SPARQL query, the query runs on the graph, and the model answers
from the result rows. The steps are:
\begin{enumerate}
  \item The prompt contains a schema summary computed once from the data:
        for each main class and each kind of career station, its properties
        with their triple counts and an example value.
  \item The prompt also gives rules and six example queries. The rules say to
        find entities by \term{rdfs:label} or \term{skos:altLabel} and never
        guess IRIs, to follow the career-station pattern or its
        \term{playedFor} shortcut, and to filter current provinces with
        \term{FILTER NOT EXISTS \{ ?p a vio:FormerProvince \}}.
  \item The generated query is parsed with rdflib, and only \term{SELECT} and
        \term{ASK} are accepted.
  \item If the query fails or returns no rows, the error is sent back to the
        model for one repair.
  \item A second prompt passes at most 30 result rows and asks for a JSON
        object with the reasoning steps and the answer. In \emph{SPARQL mode}
        this call is skipped and the query and raw rows are shown instead.
\end{enumerate}
The tab works with any OpenAI-compatible API (OpenAI, Google Gemini or a local
Ollama model), selected in \term{.env}. The examples follow one performance rule.
rdflib applies a \term{FILTER} only after joining the whole group, so the
lookup of an entity by name is written as a subquery that runs first. For the
career of Nguyễn Công Phượng this takes the query from 10.7\,s to 0.27\,s. The
tests check the repair loop, SPARQL mode, the parsing of the answer and the
rejection of write queries with a stubbed model, so they need no API key.

\subsection{Standard SPARQL Endpoint}
\label{sec:endpoint}

The built-in server offers SPARQL through the editor tab. It does not expose
the SPARQL~1.1 Protocol~\cite{feigenbaum2013protocol}, that is, a
\term{/sparql} URL that programs can call over HTTP. For that use, the dump can
be loaded into the Virtuoso triple store that is defined as an optional
service in \term{docker-compose.yml} (port 8890), next to the application
container (port 7860).
